\section{方差、协方差与双线性}


期望描述“中心位置”，方差描述“离中心有多散”。协方差则描述两个随机变量共同偏离各自中心的程度。本讲的关键是：\textbf{协方差是双线性的，所以方差可以像内积一样展开}。

\textbf{本节用到的记号：} $\operatorname{Var}(X)$ 为方差；$\operatorname{Cov}(X,Y)$ 为协方差；$\mu_X=E[X]$；$\sigma_X^2=\operatorname{Var}(X)$；$\rho(X,Y)$ 为相关系数。

\subsection{方差的定义}


\[
\operatorname{Var}(X)=E\bigl[(X-E[X])^2\bigr].
\]

它恒非负。展开平方：

\[
\begin{aligned}
\operatorname{Var}(X)
&=E[X^2-2X E[X]+(E[X])^2]\\
&=E[X^2]-2(E[X])^2+(E[X])^2\\
&=E[X^2]-(E[X])^2.
\end{aligned}
\]

这是最常用的计算公式。

\textbf{性质.}

\begin{itemize}
\item $\operatorname{Var}(aX+b)=a^2\operatorname{Var}(X)$。
\item 常数没有方差：$\operatorname{Var}(c)=0$。
\item 标准差的量纲与 $X$ 一致：$\sigma_X=\sqrt{\operatorname{Var}(X)}$。
\end{itemize}


\subsection{协方差的定义}


\[
\operatorname{Cov}(X,Y)
=E\bigl[(X-E[X])(Y-E[Y])\bigr].
\]

展开：

\[
\operatorname{Cov}(X,Y)=E[XY]-E[X]E[Y].
\]

\textbf{性质.}

\begin{itemize}
\item 对称：$\operatorname{Cov}(X,Y)=\operatorname{Cov}(Y,X)$。
\item $\operatorname{Cov}(X,X)=\operatorname{Var}(X)$。
\item 与常数无关：$\operatorname{Cov}(X,c)=0$。
\item 双线性：对任意常数 $a,b$，
\end{itemize}

\[
  \operatorname{Cov}(aX+bY,Z)
  =a\operatorname{Cov}(X,Z)+b\operatorname{Cov}(Y,Z).
\]
  对第二个变量也同理。

\textbf{相关系数.} 当 $\operatorname{Var}(X),\operatorname{Var}(Y)>0$ 时，

\[
\rho(X,Y)=\frac{\operatorname{Cov}(X,Y)}
{\sqrt{\operatorname{Var}(X)\operatorname{Var}(Y)}}.
\]

它只反映\textbf{线性相关}的程度，且 $|\rho|\le1$。

\subsection{方差的和公式}


由双线性，

\[
\begin{aligned}
\operatorname{Var}(X+Y)
&=\operatorname{Cov}(X+Y,X+Y)\\
&=\operatorname{Cov}(X,X)+\operatorname{Cov}(X,Y)
+\operatorname{Cov}(Y,X)+\operatorname{Cov}(Y,Y)\\
&=\operatorname{Var}(X)+\operatorname{Var}(Y)+2\operatorname{Cov}(X,Y).
\end{aligned}
\]

更一般地，

\[
\operatorname{Var}\Bigl(\sum_{i=1}^nX_i\Bigr)
=\sum_{i=1}^n\operatorname{Var}(X_i)
+2\sum_{1\le i<j\le n}\operatorname{Cov}(X_i,X_j).
\]

\textbf{独立情形.} 若 $X,Y$ 独立，则

\[
E[XY]=E[X]E[Y]\quad\Longrightarrow\quad
\operatorname{Cov}(X,Y)=0,
\]

于是

\[
\operatorname{Var}(X+Y)=\operatorname{Var}(X)+\operatorname{Var}(Y).
\]

\subsection{独立与不相关}


由独立可以推出

\[
E[XY]=E[X]E[Y],
\]

从而 $\operatorname{Cov}(X,Y)=0$，称 $X,Y$ \textbf{不相关}。

但反过来不成立：

\begin{quote}
\textbf{不相关（$\operatorname{Cov}=0$）不一定独立。}
\end{quote}


\textbf{反例.} 令 $X$ 在 $\{-1,0,1\}$ 上等可能，令 $Y=X^2$。则

\[
E[X]=0,\qquad E[Y]=\frac23,\qquad E[XY]=E[X^3]=0,
\]

所以 $\operatorname{Cov}(X,Y)=0$，但 $Y$ 完全由 $X$ 决定，二者显然不独立。

\textbf{记忆.} 独立 $\Rightarrow$ 不相关；不相关 $\nRightarrow$ 独立。

\subsection{协方差矩阵}


把随机向量 $\mathbf X=(X_1,\dots,X_n)^T$ 的协方差排成矩阵：

\[
\Sigma_{ij}=\operatorname{Cov}(X_i,X_j).
\]

则 $\Sigma$ 是对称半正定矩阵：对任意向量 $a$，

\[
a^T\Sigma a=\operatorname{Var}(a^T\mathbf X)\ge0.
\]

\textbf{例（Bernoulli）.} $X\sim\mathrm{Bern}(p)$：

\[
E[X]=p,\qquad E[X^2]=p,
\]

所以

\[
\operatorname{Var}(X)=p-p^2=p(1-p).
\]

\textbf{例（二项）.} 若 $X=\sum_{i=1}^nX_i$，$X_i$ 独立同分布 $\mathrm{Bern}(p)$，则

\[
\operatorname{Var}(X)=\sum_{i=1}^n\operatorname{Var}(X_i)=np(1-p).
\]

这里独立性保证交叉协方差全为 $0$。

\subsection{常见错误}


\begin{itemize}
\item 把 $\operatorname{Var}(X+Y)$ 写成 $\operatorname{Var}(X)+\operatorname{Var}(Y)$ 而漏掉 $2\operatorname{Cov}(X,Y)$。
\item 把“不相关”当成“独立”。
\item 把 $\operatorname{Cov}(X,Y)=0$ 当成“$X,Y$ 没有关系”；它只说明没有线性相关。
\item 忘记 $\operatorname{Var}(aX+b)=a^2\operatorname{Var}(X)$ 中的平方。
\end{itemize}
